\documentclass[11pt]{article}
\usepackage[utf8]{inputenc}
\usepackage[T1]{fontenc}
\usepackage{lmodern}
\usepackage[spanish]{babel}
\usepackage{amsmath,amssymb}
\usepackage{graphicx}
\usepackage[left=2cm,right=2cm,top=2cm,bottom=2cm]{geometry}
\usepackage{enumitem}
\usepackage{booktabs}
\usepackage{array}
\usepackage{xcolor}
\usepackage{framed}

\spanishdecimal{.}

\setlist[enumerate,1]{label=\textbf{\arabic*.}, leftmargin=*, itemsep=6pt}
\setlist[enumerate,2]{label=(\alph*), itemsep=4pt}
\setlength{\parindent}{0pt}
\setlength{\parskip}{4pt}
\pagestyle{plain}

\newcommand{\Ein}{E_{\text{in}}}
\newcommand{\Eout}{E_{\text{out}}}
\newcommand{\Hyp}{\mathcal{H}}
\newcommand{\dvc}{d_{\text{VC}}}
\newcommand{\vx}{\mathbf{x}}
\newcommand{\vw}{\mathbf{w}}
\newcommand{\vy}{\mathbf{y}}
\newcommand{\R}{\mathbb{R}}
\newcommand{\T}{^{\top}}
\newcommand{\E}{\mathbb{E}}
\newcommand{\Prob}{\mathbb{P}}
\newcommand{\pts}[1]{\hfill\textbf{(#1~puntos)}}
\newcommand{\linea}[1]{\rule{#1}{0.4pt}}
\newcommand{\ptsinc}[1]{{\unskip\nobreak\hfil\penalty50\hskip1em\hbox{}\nobreak\hfill(#1)\parfillskip=0pt\finalhyphendemerits=0\par}}

\newcommand{\encabezado}[2]{%
\begin{center}
{\large\bfseries Reconocimiento de Patrones}\\[2pt]
Universidad de Sonora \textperiodcentered{} Licenciatura en Ciencias de la Computación\\
Semestre 2027-1 \textperiodcentered{} Prof.\ Julio Waissman Vilanova\\[10pt]
{\Large\bfseries #1}%
\ifdefined\tipodoc\\[4pt]{\large\itshape\tipodoc}\fi
\end{center}
\ifdefined\tipodoc\else
\vspace{6pt}
\textbf{Nombre:} \hrulefill\quad \textbf{Expediente:} \rule{3.5cm}{0.4pt}\par
\vspace{6pt}
\fi
#2
\par\medskip\hrule\medskip
}

\newcommand{\espacio}[1]{\par\vspace{#1}}
\newcommand{\nuevapagina}{\newpage}


\begin{document}

\encabezado{Examen parcial \textperiodcentered{} Unidades 5 a 8}{%
Tiempo: 90 minutos. Total: 100 puntos. Responde en hojas aparte, indicando el número de problema y dejando indicadas las operaciones. Puedes usar calculadora y una hoja tamaño carta con notas escritas a mano por ti. Si te atoras en un inciso, supón el resultado y sigue: los incisos posteriores se califican con tu resultado.}

\begin{enumerate}

%% ===============================================================
\item \textbf{Un árbol de regresión y dos pasos de \textit{boosting}.} \pts{13}

Seis datos con un solo atributo:
\begin{center}
\begin{tabular}{c|cccccc}
$x$ & 1 & 2 & 3 & 4 & 5 & 6 \\ \hline
$y$ & 2 & 4 & 3 & 9 & 11 & 10
\end{tabular}
\end{center}
\begin{enumerate}
  \item Para cada umbral $t \in \{1.5, 2.5, 3.5, 4.5, 5.5\}$, calcula la suma de cuadrados de la división $x \le t$ (cada hoja predice el promedio de sus $y$). ¿Qué tocón (árbol de una división) se elige y qué predice en cada hoja? ¿Qué predice para $x = 100$? \ptsinc{4}
  \item Ahora \textit{gradient boosting} con pérdida cuadrática, tocones y tasa $\eta = 0.5$, empezando en $F_0 = \bar y$. Calcula los pseudo-residuos $r_n = y_n - F_0$ y explica por qué el primer tocón, ajustado a ellos, divide en el mismo umbral que en (a). Calcula $F_1 = F_0 + \eta\, h_1$, los residuos nuevos y el error cuadrático medio antes y después del paso. \ptsinc{5}
  \item Muestra que el segundo tocón vuelve a dividir en $t = 3.5$ (basta compararlo con $t = 4.5$) y calcula $F_2$. ¿A qué tiende $F_M$ cuando $M \to \infty$? ¿Para qué sirve entonces una $\eta$ menor que 1? \ptsinc{4}
\end{enumerate}

%% ===============================================================
\item \textbf{Promediar árboles.} \pts{12}

\begin{enumerate}
  \item Sean $Z_1, \dots, Z_B$ las predicciones de $B$ árboles en un punto $\vx$, cada una con varianza $\sigma^2$ y con correlación $\rho$ entre cada par. Demuestra que
  $\operatorname{Var}\big(\bar Z\big) = \rho\,\sigma^2 + \dfrac{1 - \rho}{B}\,\sigma^2$. \ptsinc{4}
  \item En un problema, los árboles del \textit{bagging} tienen $\sigma^2 = 100$ y $\rho = 0.5$; los de un bosque aleatorio, $\sigma^2 = 130$ y $\rho = 0.2$. Calcula la varianza del promedio de cada ensamble con $B = 1$, $B = 10$ y $B \to \infty$. ¿Por qué cada árbol del bosque es \emph{peor} y aun así el ensamble es mejor? \ptsinc{4}
  \item En el \textit{bagging} y en los bosques se usan árboles profundos; en el \textit{boosting}, árboles chicos. Explica por qué, en términos de sesgo y varianza. ¿En cuál de los dos métodos hay que elegir con cuidado el número de árboles, y por qué? \ptsinc{4}
\end{enumerate}

%% ===============================================================
\item \textbf{Las hojas de XGBoost.} \pts{12}

Ocho clientes de un banco, con su edad y si contrataron ($y = 1$) o no ($y = 0$) un depósito:
\begin{center}
\begin{tabular}{l|cccccccc}
edad & 22 & 25 & 31 & 38 & 45 & 52 & 60 & 67 \\ \hline
$y$ & 0 & 0 & 0 & 1 & 1 & 0 & 1 & 1
\end{tabular}
\end{center}
Usamos pérdida logística con $F$ en escala de log-momios, $L = \ln(1 + e^{F}) - yF$, cuyas derivadas son $g = p - y$ y $h = p(1 - p)$, con $p = 1/(1 + e^{-F})$. Se parte de $F_0 = 0$ para todos. Recuerda que una hoja con sumas $G$ y $H$ vale $w^* = -G/(H + \lambda)$ y que la ganancia de una división es
\[
\tfrac12\left[\frac{G_I^2}{H_I + \lambda} + \frac{G_D^2}{H_D + \lambda} - \frac{(G_I + G_D)^2}{H_I + H_D + \lambda}\right] - \gamma .
\]
\begin{enumerate}
  \item Calcula $g_n$ y $h_n$ para cada cliente. \ptsinc{2}
  \item Con $\lambda = 1$ y $\gamma = 0$, calcula $G$, $H$ y $w^*$ de cada hoja, y la ganancia, para las divisiones edad~$\le 34.5$ y edad~$\le 48.5$. ¿Cuál se elige? \ptsinc{6}
  \item ¿Con qué valores de $\gamma$ la raíz ya no se divide? \ptsinc{2}
  \item Con tasa de aprendizaje $\eta = 0.3$ y la división elegida, ¿qué probabilidad de contratar le asigna el modelo, después de este árbol, a un cliente de 25 años y a uno de 60? \ptsinc{2}
\end{enumerate}

%% ===============================================================
\item \textbf{Retropropagación con salida logística.} \pts{15}

Una red con entrada $\vx \in \R^2$, dos neuronas ocultas ReLU (sin sesgo) y una salida logística:
\[
\mathbf z^{(1)} = W^{(1)}\vx, \quad \mathbf a^{(1)} = \operatorname{ReLU}\big(\mathbf z^{(1)}\big), \quad z^{(2)} = \mathbf w^{(2)\top}\mathbf a^{(1)} + b^{(2)}, \quad p = \sigma\big(z^{(2)}\big),
\]
con $W^{(1)} = \begin{pmatrix} 1 & 0 \\ 1 & 2 \end{pmatrix}$, $\mathbf w^{(2)} = (2, -1)$, $b^{(2)} = -1$ y entropía cruzada $E = -y\ln p - (1 - y)\ln(1 - p)$. El ejemplo es $\vx = (1, -1)$ con $y = 0$.
\begin{enumerate}
  \item Calcula $\mathbf z^{(1)}$, $\mathbf a^{(1)}$, $z^{(2)}$, $p$ y $E$. \ptsinc{3}
  \item Demuestra que $\delta^{(2)} = \partial E/\partial z^{(2)} = p - y$. (Usa $\sigma'(z) = \sigma(z)(1 - \sigma(z))$). \ptsinc{3}
  \item Calcula $\partial E/\partial \mathbf w^{(2)}$, $\partial E/\partial b^{(2)}$, $\boldsymbol\delta^{(1)}$ y $\partial E/\partial W^{(1)}$. \ptsinc{4}
  \item ¿Qué pasaría si todos los pesos y sesgos se inicializaran en cero? Y si se inicializaran todos con el mismo valor $c \ne 0$, ¿qué tendrían en común las neuronas ocultas durante todo el entrenamiento? \ptsinc{3}
  \item Para verificar tu gradiente con diferencias finitas, ¿usarías $\frac{E(\theta + h) - E(\theta)}{h}$ o $\frac{E(\theta + h) - E(\theta - h)}{2h}$? ¿Por qué? \ptsinc{2}
\end{enumerate}

%% ===============================================================
\item \textbf{Convoluciones y atención.} \pts{15}

Una red para imágenes de satélite de $3 \times 64 \times 64$:
\begin{center}
\begin{tabular}{ll}
\toprule
capa & detalle \\ \midrule
conv 1 & 16 filtros de $3 \times 3$, relleno 1, paso 1, ReLU \\
\textit{max-pooling} & $2 \times 2$, paso 2 \\
conv 2 & 32 filtros de $3 \times 3$, relleno 1, paso 1, ReLU \\
\textit{max-pooling} & $2 \times 2$, paso 2 \\
conv 3 & 64 filtros de $3 \times 3$, sin relleno, paso 2, ReLU \\
promedio global & un promedio por canal \\
lineal & 10 clases \\
\bottomrule
\end{tabular}
\end{center}
\begin{enumerate}
  \item Escribe el tamaño de la salida y el número de parámetros de cada capa, y el total. \ptsinc{6}
  \item ¿De qué tamaño es el campo receptivo de una neurona de conv 3? \ptsinc{3}
  \item En una capa de autoatención con $d_k = 4$, la consulta de la posición 2 es $\mathbf q = (1, 0, 1, 2)$. Las claves y valores de las posiciones 1, 2 y 3 son
  \[
  \mathbf k_1 = (1, 1, 1, 1),\ \ \mathbf k_2 = (2, 0, 0, 1),\ \ \mathbf k_3 = (0, 1, -1, 0); \qquad \mathbf v_1 = (3, 0),\ \ \mathbf v_2 = (0, 3),\ \ \mathbf v_3 = (1, 1).
  \]
  Calcula los pesos de atención escalados y la salida de la posición 2 sin máscara, y luego con máscara causal. \ptsinc{4}
  \item ¿Cómo crece el costo de la autoatención con la longitud $T$ de la secuencia? ¿Qué pasa con el costo si se duplica el contexto? \ptsinc{2}
\end{enumerate}

%% ===============================================================
\item \textbf{La temperatura de mañana.} \pts{17}

La temperatura media diaria de Hermosillo es $y_t = c(d_t) + a_t$, donde $c(d)$ es la climatología (lo normal para el día del año $d$) y las anomalías siguen un AR(1), $a_t = 0.8\, a_{t-1} + \varepsilon_t$, con $\varepsilon_t$ ruido blanco gaussiano de desviación estándar $\sigma = 1.7$\,°C.
\begin{enumerate}
  \item ¿Cuál es la desviación estándar de las anomalías? \ptsinc{2}
  \item Hoy la temperatura media fue de 29.0\,°C y lo normal era 25.0\,°C. Lo normal para los próximos tres días es 25.2, 25.4 y 25.6\,°C. Pronostica los tres días, da el intervalo del 95\,\% para mañana, y di a qué tiende el intervalo cuando el horizonte crece. \ptsinc{5}
  \item Un compañero quiere pronosticar la temperatura media con 7 días de anticipación usando un bosque aleatorio. Para pronosticar el día $t$ usa como atributos: la temperatura de los días $t - 7$ y $t - 8$, la media de los 30 días anteriores a $t$, la humedad y el viento medidos el día $t$, y el día del año. Lo evalúa con validación cruzada de 5 pliegues al azar y reporta un MAE de 1.1\,°C. Señala tres problemas y describe cómo lo evaluarías correctamente. \ptsinc{5}
  \item Para dar intervalos con un modelo de aprendizaje se usa la pérdida cuantil
  \[ \rho_\tau(u) = u\,\big(\tau - \mathbf 1[u < 0]\big), \qquad u = y - q . \]
  Demuestra que, si $Y$ tiene densidad y función de distribución $F$, el $q$ que minimiza $\E\big[\rho_\tau(Y - q)\big]$ cumple $F(q) = \tau$. ¿Cómo armarías un intervalo del 80\,\% y cómo comprobarías que lo es? \ptsinc{5}
\end{enumerate}

%% ===============================================================
\item \textbf{Un año en producción.} \pts{16}

Un modelo de rentas de bicicletas por hora, entrenado con 2011, está en servicio desde el 1 de enero de 2012 y registra en su bitácora cada entrada y cada predicción.
\begin{enumerate}
  \item Las predicciones de 2011 se repartieron en cuatro cubetas con proporciones $p = (0.4, 0.3, 0.2, 0.1)$; las de la primera semana de febrero de 2012, $q = (0.3, 0.3, 0.2, 0.2)$. Calcula el PSI. \ptsinc{3}
  \item Sin ningún cambio, $\E[\text{PSI}] \approx (K - 1)\big(\frac1N + \frac1M\big)$. Si se compara una semana (168 horas) contra otra semana, ¿cuánto vale ese piso con $K = 4$ y con $K = 10$ cubetas? ¿Qué le pasa a la regla «PSI~$> 0.1$ es cambio» con $K = 10$? \ptsinc{3}
  \item El tablero compara la temperatura de cada mes de 2012 contra \emph{todo} 2011 y marca alarma todos los meses (PSI entre 1 y 5). ¿Por qué? ¿Qué referencia usarías? \ptsinc{3}
  \item Con 100,000 predicciones por lado, la prueba de Kolmogorov–Smirnov da $D = 0.01$. Compara con el umbral del 5\,\%, $1.358\sqrt{(n+m)/(nm)}$. ¿Hay que reentrenar? \ptsinc{3}
  \item En octubre llegan las etiquetas: la RECM pasó de 47 (validación en 2011) a 125, y el modelo subestima todos los meses, mientras los PSI de las predicciones siguen bajos. ¿Qué tipo de cambio es? Describe cómo pondrías en producción un modelo reentrenado usando un registro de modelos, y qué tendrías que decidir sobre los datos de reentrenamiento. \ptsinc{4}
\end{enumerate}

\end{enumerate}

\end{document}
